% Chapter 2, Section 2 _Linear Algebra_ Jim Hefferon
%  http://joshua.smcvt.edu/linearalgebra
%  2001-Jun-10
\section{线性无关}
前一节展示了如何将一个向量空间
理解为一个张成，
即理解为它的某些元素的毫无限制的线性组合。
例如，次数不超过一次的多项式构成的空间 $\set{a+bx\suchthat a,b\in\Re}$
由集合 $\set{1,x}$ 张成。
前一节还表明，一个空间可以有许多张成它的集合。
张成该多项式空间的另外两个集合是
$\set{1,2x}$ 与 $\set{1,x,2x}$。

在那一节的末尾，我们把一些张成集描述为“最小的”，
但从未精确定义过这个词。
% We could mean one of two things.
一种含义是：一个张成集是最小的，若它
在所有具有相同张成结果的集合中含有个数最少的成员，
这样 $\set{1,x,2x}$ 就不是最小的，因为它有
三个成员，而我们能给出张成同一空间的两元素集合。
或者，我们也可以认为，当一个张成集不含任何
去掉后不改变张成结果的元素时，它是最小的。
按这种含义，$\set{1,x,2x}$ 不是最小的，因为
去掉 \( 2x \) 而得到 \( \set{1,x} \)，张成结果
并不改变。

第一种最小性的含义看来是一个整体性的要求，
因为要检验一个张成集是否最小，
我们似乎必须考察所有张成该空间的集合，
并找出其中元素个数最少的一个。
第二种含义则是局部性的，因为
我们只需考察这个集合本身，考虑
去掉或保留各种元素时的张成结果。
例如，按照第二种含义，
我们可以将 $\set{1,x,2x}$ 的张成
与 $\set{1,x}$ 的张成相比较，并注意到
$2x$ 是一个“重复”，因为
去掉它并不会使张成缩小。

本节将采用“最小张成集”的第二种含义，
原因在于它在技术上较为方便。
然而，本书最重要的结果是这两种含义
是一致的。
我们将在下一节证明这一点。

  
 


 


\subsection{定义与例题}

我们在第一章中见过“重复”。
在那里，高斯消元法把它们化成了
$0=0$ 型的方程。

\begin{example}  \label{ex:StaticsLIAndLD}
回忆第一章开头的 \hyperlink{ex:Statics}{静力学}
例题。
我们由那对质量未知的物体得到了两次平衡，一次
在 \( 40 \)~厘米与 \( 15 \)~厘米处，另一次
在 \( -50 \)~厘米与 \( 25 \)~厘米处，
然后我们算出了那些质量值。
假如第二次平衡得到的是
在 \( 20 \)~厘米与 \( 7.5 \)~厘米处，
那么对由此得到的二方程、二未知量的方程组使用高斯消元法
将不会给出解，而是给出一个 $0=0$ 的方程
连同一个含自由变量的方程。
直观地说，问题在于 \( \rowvec{20 &7.5} \) 是
$\rowvec{40 &15}$ 的一半，也就是说，
$\rowvec{20 &7.5}$ 位于集合 $\set{\rowvec{40 &15}}$ 的张成之中，
因而是重复的数据。
那样我们就是在试图
用本质上仅有一条的信息去求解一个含两个未知量的问题。
\end{example}
  
若 $\vec{v}\in\spanof{S}$，从而它依赖于集合~$S$
的元素——也就是说，它能用该集合的元素表示出来，
$\vec{v}=c_1\vec{s}_1+\cdots+c_n\vec{s}_n$——
我们就把 $\vec{v}$ 取作集合~$S$ 中各向量的一个“重复”。

\begin{lemma} \label{lm:AddVecSpanNoGrowIffVecInSpan}
%<*lm:AddVecSpanNoGrowIffVecInSpan>
设 $V$ 是一个向量空间，$S$ 是该空间的一个子集，而 $\vec{v}$
是该空间的一个元素，则
$\spanof{S\union\set{\vec{v}}} = \spanof{S}$
当且仅当
$\vec{v}\in\spanof{S}$。
%</lm:AddVecSpanNoGrowIffVecInSpan>
\end{lemma}

\begin{proof}
%<*pf:AddVecSpanNoGrowIffVecInSpan0>
“当且仅当”的一半是直接的：若 $\vec{v}\notin\spanof{S}$，则
这两个集合不相等，因为 $\vec{v}\in\spanof{S\union\set{\vec{v}}}$。

对于另一半，设 $\vec{v}\in\spanof{S}$，
于是对某些标量~$c_i$ 及向量 $\vec{s}_i\in S$，
有 $\vec{v}=c_1\vec{s}_1+\cdots+c_n\vec{s}_n$。
我们将用相互包含来证明集合
$\spanof{S\union\set{\vec{v}}}$ 与 $\spanof{S}$ 相等。
包含关系 $\spanof{S\union\set{\vec{v}}}\supseteq\spanof{S}$ 是显然的。
%</pf:AddVecSpanNoGrowIffVecInSpan0>

%<*pf:AddVecSpanNoGrowIffVecInSpan1>
为证明另一方向的包含，设 $\vec{w}$ 是~$\spanof{S\union\set{\vec{v}}}$
中的一个元素。
那么 $\vec{w}$ 是 $S\union\set{\vec{v}}$ 中元素的线性组合，
我们可以把它写成
$\vec{w}=c_{n+1}\vec{s}_{n+1}+\cdots+c_{n+k}\vec{s}_{n+k}+c_{n+k+1}\vec{v}$。
（$\vec{w}$ 的等式中的一些 $\vec{s}_i$ 有可能与
$\vec{v}$ 的等式中的某些相同，但这无关紧要。）
把 $\vec{v}$ 展开代入。
\begin{equation*}
  \vec{w}=c_{n+1}\vec{s}_{n+1}+\cdots+c_{n+k}\vec{s}_{n+k}
            +c_{n+k+1}\cdot(c_1\vec{s}_1+\cdots+c_n\vec{s}_n)
\end{equation*}
把右端看成来自~$S$ 的向量的线性组合的
线性组合。
因此 $\vec{w}\in\spanof{S}$。
%</pf:AddVecSpanNoGrowIffVecInSpan1>
\end{proof}

本节开头的讨论涉及的是去掉向量，而不是
添加向量。

\begin{corollary}  \label{cor:OmitVecNotShrinkSpanIffVecIsDep}
%<*cor:OmitVecNotShrinkSpanIffVecIsDep>
对于 $\vec{v}\in S$，
去掉该向量不会使张成缩小
当且仅当该向量依赖于集合中的
其他向量。
也就是说，
$\spanof{S}=\spanof{S-\set{\vec{v}}}$
当且仅当 $\vec{v}\in\spanof{S-\set{\vec{v}}}$。
%</cor:OmitVecNotShrinkSpanIffVecIsDep>
\end{corollary}

因此，要知道去掉一个向量是否会减小
张成，就需要知道该向量是否是集合中
其他向量的线性组合。

\begin{definition}\label{def:LinInd}
%<*df:LinInd>
在任何向量空间中，一个由向量组成的集合
称为\definend{线性无关的}\index{linearly independent}%
\index{sets!dependent, independent}
若集合中没有任何元素是集合中
其他元素的线性组合。\footnote{另见\nearbyremark{rem:WhyLIUsesMultiset}。}\spacefactor=1000\ 
否则，该集合是
\definend{线性相关的}\index{linearly dependent}。
%</df:LinInd>
\end{definition}
% A multiset subset of a vector space is
% \definend{linearly independent}\index{linearly independent}%
% \index{sets!dependent, independent}
% if none of its elements is a linear combination of the 
% others.\footnote{More information on multisets is in the appendix. 
% Most of the time we won't need the 
% set-multiset distinction and we will
% follow the standard terminology of referring to a linearly independent or 
% dependent `set'.
% \nearbyremark{rem:WhyLIUsesMultiset} explains why 
% the definition requires a multiset.}\spacefactor=1000\ 
% Otherwise it is \definend{linearly dependent}\index{linearly dependent}.

因此，集合 $\set{\vec{s}_0,\ldots,\vec{s}_n}$ 是线性无关的
如果不存在等式
$\vec{s}_i=c_0\vec{s}_0+\ldots+c_{i-1}\vec{s}_{i-1}+c_{i+1}\vec{s}_{i+1}+\ldots+c_n\vec{s}_n$。
定义中“其他”一词的用法意味着，
用 $\vec{s}_i=1\cdot\vec{s}_i$ 把 $\vec{s}_i$ 写成线性组合
是不算数的。

%<*LinInd>
注意到，
尽管把一个向量写成其他向量的组合的这种写法
\begin{equation*}
   \vec{s}_0=\lincombo{c}{\vec{s}}
\end{equation*}
在视觉上突出了 \( \vec{s}_0 \)，但在代数上
该向量在这个等式中并无任何特殊之处。
对于任何系数 $c_i$ 非~$0$ 的 \( \vec{s}_i \)，
我们可以改写以解出 \( \vec{s}_i \)。
\begin{equation*}
   \vec{s}_i=(1/c_i)\vec{s}_0+\dots
              +(-c_{i-1}/c_i)\vec{s}_{i-1}+(-c_{i+1}/c_i)\vec{s}_{i+1}
              +\dots+(-c_n/c_i)\vec{s}_n
\end{equation*}
当我们不想突出任何向量时，
我们会改而说
\( \vec{s}_0,\vec{s}_1,\dots,\vec{s}_n \) 处于一个
\definend{线性关系}\index{linear relationship}%
\index{relationship!linear} 
之中，并把所有向量放到等号的同一侧。
%</LinInd>
下面的结论以这种方式重新表述了线性无关的定义。
这通常是我们判定
一个有限集相关还是无关时所用的方法。

\begin{lemma}   \label{le:LDIffANonTrivLinRel}
%<*lm:LDIffANonTrivLinRel>
向量空间的一个子集 \( S \) 是线性无关的，当且仅当
在其元素之中
唯一的线性关系
$ c_1\vec{s}_1+\dots+c_n\vec{s}_n=\zero$
是平凡的，\( c_1=0,\dots,\,c_n=0 \)
（其中当 $i\neq j$ 时 $\vec{s}_i\neq\vec{s}_j$）。
%</lm:LDIffANonTrivLinRel>
\end{lemma}

\begin{proof}
%<*pf:LDIffANonTrivLinRel>
若 \( S \) 是线性无关的，则没有任何向量
$\vec{s}_i$
是 $S$ 中其他向量的线性组合，
因而不存在这样的线性关系，其中某些
$\vec{s}\,$ 的系数非零。

若 \( S \) 不是线性无关的，则某个 \( \vec{s}_i \) 是 $S$ 中
其他向量的线性组合
$\vec{s}_i=c_1\vec{s}_1+\dots+c_{i-1}\vec{s}_{i-1}
    +c_{i+1}\vec{s}_{i+1}+\dots+c_n\vec{s}_n$
。
从两边减去 $\vec{s}_i$
得到一个关系
其中含有非零系数，
即 \( \vec{s}_i \) 前面的 \( -1 \)。
%</pf:LDIffANonTrivLinRel>
\end{proof}

\begin{example}  \label{ex:StaticsLI}
在宽度为二的行向量构成的向量空间中，二元素集合
\( \set{ \rowvec{40 &15},\rowvec{-50 &25}} \) 是线性无关的。
为验证这一点，取
\begin{equation*}
  c_1\cdot\rowvec{40 &15}+c_2\cdot\rowvec{-50 &25}=\rowvec{0 &0}
\end{equation*}
并求解由此得到的方程组。
\begin{equation*}
  \begin{linsys}{2}
    40c_1  &-  &50c_2  &=  &0  \\
    15c_1  &+  &25c_2  &=  &0  
   \end{linsys}
  \grstep{-(15/40)\rho_1+\rho_2}
  \begin{linsys}{2}
     40c_1  &- &50c_2       &=  &0  \\
            &  &(175/4)c_2  &=  &0  
   \end{linsys}
\end{equation*}
\( c_1 \) 与 \( c_2 \) 都为零。
因此，所给的两个行向量之间唯一的线性关系
是平凡的关系。

在同一向量空间中，集合
\( \set{ \rowvec{40 &15},\rowvec{20 &7.5}} \) 是线性相关的，因为
我们可以满足
% \begin{equation*}
$c_1\cdot \rowvec{40 &15}+c_2\cdot\rowvec{20 &7.5}=\rowvec{0 &0}$
% \end{equation*}
其中 \( c_1=1 \) 且 \( c_2=-2 \)。
\end{example}

\begin{example}
集合 \( \set{1+x,1-x} \) 在 \( \polyspace_2 \) 中是线性无关的，
该空间由实系数的、次数不超过二的多项式组成，因为
\begin{equation*}
   0+0x+0x^2
   =
   c_1(1+x)+c_2(1-x)
   =
   (c_1+c_2)+(c_1-c_2)x+0x^2
\end{equation*}
给出
\begin{equation*}
  \begin{linsys}{2}
    c_1  &+  &c_2  &=  &0  \\
    c_1  &-  &c_2  &=  &0  
   \end{linsys}
  \grstep{-\rho_1+\rho_2}
  \begin{linsys}{2}
     c_1  &+  &c_2  &=  &0  \\
          &   &2c_2 &=  &0
  \end{linsys}
\end{equation*}
这是因为，多项式相等当且仅当它们的系数相等。
因此，$\polyspace_2$ 中这两个成员之间的
唯一线性关系是平凡的。
\end{example}

\begin{remark}
引理之所以指明 $\vec{s}_i\neq\vec{s}_j$
（当 $i\neq j$ 时），是因为显然，若某个向量~$\vec{s}$ 出现
两次，我们就能得到一个非平凡的
$c_1\vec{s}_1+\dots+c_n\vec{s}_n=\zero$，
只需取相应的系数为 $1$ 和~$-1$。
此外，若某个向量在表达式中出现不止一次，则
我们总可以把系数合并起来。

注意，该引理允许
与出现不止一次相反的情形，即~$S$ 中的
某些向量根本不出现。
例如，若~$S$ 是无限集，那么由于
线性关系只涉及有限多个向量，
任何一个这样的关系都会略去 $S$ 的许多向量。
不过，还要注意，若~$S$ 是有限集，则在方便时
我们可以取一个组合
$c_1\vec{s}_1+\dots+c_n\vec{s}_n$，使 $S$ 的每个向量恰好出现
一次。
若有向量缺失，我们可以用系数~$0$
把它补上。
\end{remark}

\begin{example}
这一矩阵
\begin{equation*}
  A=
  \begin{mat}[r]
    2  &3  &1  &0  \\
    0  &-1 &0  &-2 \\
    0  &0  &0  &1
  \end{mat}
\end{equation*}
的各行构成一个线性无关的集合。
这一点在本例中容易验证，而且还可以回忆
引理~One.III.\ref{le:EchFormNoLinCombo}
表明，任何阶梯形矩阵的各行都构成
一个线性无关的集合。
\end{example}

\begin{example}
在 \( \Re^3 \) 中，设
\begin{equation*}
   \vec{v}_1=\colvec{3 \\ 4 \\ 5}
   \quad
   \vec{v}_2=\colvec{2 \\ 9 \\ 2}
   \quad
   \vec{v}_3=\colvec{4 \\ 18 \\ 4}
\end{equation*}
则集合 \( S=\set{\vec{v}_1,\vec{v}_2,\vec{v}_3} \)
是线性相关的，因为下面是一个关系
\begin{equation*}
  0\cdot\vec{v}_1
  +2\cdot\vec{v}_2
  -1\cdot\vec{v}_3
  =\zero
\end{equation*}
其中并非所有标量都为零
（某些标量为零
这一事实无关紧要）。
\end{example}

% \begin{remark}  \label{rem:WhyLIIsNonTrivLinRel}
上面的例题说明了为什么，
尽管 \nearbydefinition{def:LinInd} 对无关的含义是更清晰的
陈述，
\nearbylemma{le:LDIffANonTrivLinRel} 却更适合
计算。
直接从定义出发，想要计算 $S$ 是否
线性无关的人会先设
\( \vec{v}_1=c_2\vec{v}_2+c_3\vec{v}_3 \)
并得出不存在这样的 $c_2$ 与 $c_3$ 的结论。
但只知道第一个向量不
依赖于另外两个向量是不够的。
此人还必须继续尝试
\( \vec{v}_2=c_1\vec{v}_1+c_3\vec{v}_3 \)，以便
找到相关关系 $c_1=0$、\( c_3=1/2 \)。
\nearbylemma{le:LDIffANonTrivLinRel}
只用一次计算便得到同样的结论。
% \end{remark}

\begin{example} \label{ex:EmSetLI}
向量空间的空子集\index{sets!empty}是线性无关的。
在它的成员之间没有非平凡的线性关系，因为它
没有成员。
\end{example}

\begin{example} \label{ex:SetWithZeroVecLD}
在任何向量空间中，任何包含零向量的子集都是线性
相关的。
一个例子是，在二次多项式（次数不超过二）空间 $\polyspace_2$ 中，
子集 $\set{1+x,x+x^2,0}$。
它是线性
相关的，因为
$0\cdot\vec{v}_1+0\cdot\vec{v}_2+1\cdot\zero=\zero$ 是一个非平凡的
关系，因为并非所有系数都为零。
\end{example}

有一个微妙之处，我们将多次见到它，而且它
与前面的例题有关。
它是关于平凡和的，即空集的和。
看如何定义平凡和的一种方式
是考虑如下递进：
$\vec{v}_1+\vec{v}_2+\vec{v}_3$，接着是
$\vec{v}_1+\vec{v}_2$，接着是
$\vec{v}_1$。
三个向量的和与两个向量的和之差
是 $\vec{v}_3$。
而两个向量的和与一个向量的和之差
是 $\vec{v}_2$。
再进一步过渡到平凡和——零个向量的和——时，
我们可以预期减去~$\vec{v}_1$。
于是我们定义零个向量的和为零向量。

与前面例题的联系是：若零向量在一个集合中，
则该集合有一个元素
是该集合中其他向量组成的某个子集的组合，具体地说，
零向量是空子集的组合。
即使集合 $S=\set{\zero}$ 也是线性相关的，因为 $\zero$ 是
空集的和，而空集是~$S$ 的子集。

\begin{remark} \label{rem:WhyLIUsesMultiset} %\cite{Velleman} 
线性无关的定义，
\nearbydefinition{def:LinInd}，提到的是由向量组成的“集合”。
集合是我们最熟悉的一类汇集方式，而且
实际上每个人在这种语境下都使用“集合”一词。
但为完整起见，我们要指出，就这一目的而言，集合这种汇集
并不十分合适。

回忆一下，集合是具有两条性质的汇集：
(i)~次序无关紧要，
于是集合 $\set{1,2}$ 等于集合 $\set{2,1}$；以及
(ii)~重复元素合并，于是集合 $\set{1,1,2}$ 等于集合
$\set{1,2}$。

现在考虑这一矩阵的化简。
\begin{equation*}
  \begin{mat}
    1 &1 &1 \\
    2 &2 &2 \\
    1 &2 &3
  \end{mat}
  \grstep{(1/2)\rho_2}
  \begin{mat}
    1 &1 &1 \\
    1 &1 &1 \\
    1 &2 &3
  \end{mat}
\end{equation*}
在左边，矩阵行的集合
$\set{\rowvec{1 &1 &1}, \rowvec{2 &2 &2}, \rowvec{1 &2 &3}}$
是线性相关的。
在右边，行的集合是
$\set{\rowvec{1 &1 &1}, \rowvec{1 &1 &1}, \rowvec{1 &2 &3}}$。
由于重复元素合并，它等于
集合
$\set{\rowvec{1 &1 &1}, \rowvec{1 &2 &3}}$，
而它是线性无关的。
这就有问题了，因为高斯消元法理应保持线性相关性。

也就是说，严格来讲，
我们需要一类重复元素不合并的汇集。
次序无关紧要且重复元素不合并的汇集
称为\definend{多重集}。\index{sets!multiset}%
\index{multiset}\index{linear independence!multiset}

然而，尽管坚持完全正确有好处，
背离“集合”这一标准术语也会有其
自身的陷阱，所以我们将继续使用这个词。
以后，我们偶尔需要在不让重复元素合并的情况下取组合，
并且我们将这样做而不
再加说明。
\end{remark}

\begin{corollary} \label{cor:LiSetMinSpanSet}
%<*cor:LiSetMinSpanSet>
集合 $S$ 是线性无关的，当且仅当
对任何 $\vec{v}\in S$，去掉它都会使张成缩小：
$\spanof{S-\set{v}}\subsetneq\spanof{S}$。
%</cor:LiSetMinSpanSet>
\end{corollary}

\begin{proof}
%<*pf:LiSetMinSpanSet>
这由 \nearbycorollary{cor:OmitVecNotShrinkSpanIffVecIsDep} 得到。
若 $S$ 是线性无关的，则它的任何向量都不依赖于
其他元素，因此去掉任何向量都会使张成缩小。
若 $S$ 不是线性无关的，则它包含一个依赖于
集合中其他元素的向量，而去掉该向量
不会使张成缩小。
%</pf:LiSetMinSpanSet>
\end{proof}

因此，一个张成集是最小的，当且仅当它是线性无关的。

前面的结论处理的是从线性无关集合中去掉元素。
下面的结论则添加元素。

\begin{lemma}  \label{lm:AddVecLiSetIsLiIffVecNotInSpan}
%<*lm:AddVecLiSetIsLiIffVecNotInSpan>
设 $S$ 是线性无关的，且 $\vec{v}\notin S$。
那么，
集合 $S\union\set{\vec{v}}$ 是线性无关的，当且仅当
$\vec{v}\notin\spanof{S}$。
%</lm:AddVecLiSetIsLiIffVecNotInSpan>
\end{lemma}

\begin{proof}
%<*pf:AddVecLiSetIsLiIffVecNotInSpan0>
% Assume that $S$ is linearly independent and that $\vec{v}\notin S$.
我们将证明 $S\union\set{\vec{v}}$ 不是线性无关的
当且仅当 $\vec{v}\in\spanof{S}$。

先设 $\vec{v}\in\spanof{S}$。
把 $\vec{v}$ 表示为组合 $\vec{v}=c_1\vec{s}_1+\cdots+c_n\vec{s}_n$。
将其改写为 $\zero=c_1\vec{s}_1+\cdots+c_n\vec{s}_n-1\cdot\vec{v}$。
由于 $\vec{v}\notin S$，它不等于任何 $\vec{s}_i$，所以这是
$S\union\set{\vec{v}}$ 的元素之间的一个非平凡线性相关关系。
因此该集合不是线性无关的。
%</pf:AddVecLiSetIsLiIffVecNotInSpan0>

%<*pf:AddVecLiSetIsLiIffVecNotInSpan1>
现在设 $S\union\set{\vec{v}}$ 不是线性无关的，并
考虑其成员之间的一个非平凡相关关系
$\zero=c_1\vec{s}_1+\cdots+c_n\vec{s}_n+c_{n+1}\cdot\vec{v}$。
若 $c_{n+1}=0$，则那是~$S$ 的元素之间的一个相关关系，但
我们假设了~$S$ 是无关的，所以 $c_{n+1}\neq 0$。
将方程改写
为 $\vec{v}=(c_1/c_{n+1})\vec{s}_1+\cdots+(c_n/c_{n+1})\vec{s}_n$
即得 $\vec{v}\in\spanof{S}$
%</pf:AddVecLiSetIsLiIffVecNotInSpan1>
\end{proof}



\begin{example} \label{ex:LinindSetsAndSuper}
$\Re^3$ 的这个子集是线性无关的。
\begin{equation*}
  S
   =\set{\colvec{1 \\ 0 \\ 0}}
\end{equation*}
$S$ 的张成是 $x$ 轴。
这里有两个扩集，一个是线性相关的，另一个
是线性无关的。
\begin{center}
     相关：
     \( \set{
         \colvec{1 \\ 0 \\ 0},
         \colvec{-3 \\ 0 \\ 0}} \)      
     \qquad
     无关：
     \( \set{
         \colvec{1 \\ 0 \\ 0},
         \colvec{0 \\ 1 \\ 0}} \)      
\end{center}
我们得到
相关扩集的方式是添加一个来自 $x$ 轴的向量，
因此张成没有扩大。
我们得到无关扩集的方式是添加一个
不在 $\spanof{S}$ 中的向量，因为它的 $y$~分量非零，
这使张成扩大了。

对于无关集
\begin{equation*}
  S
   =\set{\colvec{1 \\ 0 \\ 0},
          \colvec{0 \\ 1 \\ 0}
                 }
\end{equation*}
张成 \( \spanof{S} \) 是 \( xy \) 平面。
这里有两个扩集。
\begin{center}
     相关：
     \( \set{
         \colvec{1 \\ 0 \\ 0},
         \colvec{0 \\ 1 \\ 0},
         \colvec{3 \\ -2 \\ 0} } \)       
     \qquad
     无关：
     \( \set{
         \colvec{1 \\ 0 \\ 0},
         \colvec{0 \\ 1 \\ 0},
         \colvec{0 \\ 0 \\ 1} } \)    
\end{center}
如上，相关扩集中
添加的成员来自 $\spanof{S}$，即 \( xy \) 平面，而
无关扩集中
添加的成员则来自该张成之外。

最后，考虑这个无关集
\begin{equation*}
   S =\set{\colvec{1 \\ 0 \\ 0},
          \colvec{0 \\ 1 \\ 0},
          \colvec{0 \\ 0 \\ 1}        } 
\end{equation*}
且 \( \spanof{S}=\Re^3 \)。
我们能得到一个线性相关的扩集。
\begin{center}
     相关：
     \( \set{
         \colvec{1 \\ 0 \\ 0},
         \colvec{0 \\ 1 \\ 0},
         \colvec{0 \\ 0 \\ 1},
         \colvec{2 \\ -1 \\ 3} } \)     
\end{center}
但 $S$ 没有线性无关的扩集。
看到这一点的一种方式是注意
对于任何要添加到 $S$ 中的向量，方程
\begin{equation*}
  \colvec{x \\ y \\ z}
  =c_1\colvec{1 \\ 0 \\ 0}
   +c_2\colvec{0 \\ 1 \\ 0}
   +c_3\colvec{0 \\ 0 \\ 1}
\end{equation*}
都有解 $c_1=x$、$c_2=y$、$c_3=z$。
看到这一点的另一种方式是：
我们无法添加任何来自张成 $\spanof{S}$ 之外的向量，因为这个
张成就是 $\Re^3$。
\end{example}

\begin{corollary} \label{th:AlwaysAnLDSubset}
%<*th:AlwaysAnLDSubset>
在向量空间中，
任何有限集都有具有相同张成的线性无关子集。
%</th:AlwaysAnLDSubset>
\end{corollary}

\begin{proof}
%<*pf:AlwaysAnLDSubset0>
若 \( S=\set{ \vec{s}_1,\dots,\vec{s}_n} \) 是线性无关的，
那么 $S$ 本身就满足结论，所以
设它是线性相关的。
%</pf:AlwaysAnLDSubset0>

%<*pf:AlwaysAnLDSubset1>
由相关的定义，$S$ 包含
一个向量 $\vec{v}_1$，它是
其他向量的线性组合。
定义集合 \( S_1=S-\set{\vec{v}_1} \)。
由 \nearbycorollary{cor:OmitVecNotShrinkSpanIffVecIsDep} % {lm:ShrinkSpanByRemovingNonRepeat} 
张成不缩小 \( \spanof{S_1}=\spanof{S} \)。
%</pf:AlwaysAnLDSubset1>

%<*pf:AlwaysAnLDSubset2>
若 \( S_1 \) 是线性无关的，那么就完成了。
否则迭代：
取一个向量 $\vec{v}_2$，
它是 $S_1$ 的
其他成员的线性组合，将其舍弃，
得到 \( S_2=S_1-\set{\vec{v}_2} \)，
满足 \( \spanof{S_2}=\spanof{S_1} \)。
重复这一过程，直到出现线性无关集 $S_j$；
它必定最终会出现，因为 \( S \) 是有限的，
而空集是线性无关的。
%</pf:AlwaysAnLDSubset2>
（严格地讲，这一论证使用
关于 $S$ 中元素个数的归纳法。
\nearbyexercise{exer:FillIndDetProofSetHasLISub} 要求给出细节。）
\end{proof}

因此，如果我们有一个线性相关的集合，那么我们可以在不改变
张成的情况下，通过舍弃我们所说的“重复”向量
来缩减它。

\begin{example}  \label{ex:ShrinkSetSameSpan}
这个集合张成 \( \Re^3 \)（验证是常规的），
但它不是线性无关的。
\begin{equation*}
  S=\set{\colvec[r]{1 \\ 0 \\ 0},
     \colvec[r]{0 \\ 2 \\ 0},
     \colvec[r]{1 \\ 2 \\ 0},
     \colvec[r]{0 \\ -1 \\ 1},
     \colvec[r]{3 \\ 3 \\ 0}  }
\end{equation*}
我们将计算应去掉哪些向量，
以得到一个无关但
具有相同张成的子集。
这一线性关系
\begin{equation*}
  c_1\colvec[r]{1 \\ 0 \\ 0}
  +c_2\colvec[r]{0 \\ 2 \\ 0}
  +c_3\colvec[r]{1 \\ 2 \\ 0}
  +c_4\colvec[r]{0 \\ -1 \\ 1}
  +c_5\colvec[r]{3 \\ 3 \\ 0}
  =\colvec[r]{0 \\ 0 \\ 0}
  \qquad\tag{$*$}
\end{equation*}
给出方程组
\begin{equation*}
  \begin{linsys}{5}
     c_1  &   &      &+  &c_3   &+  &    &+  &3c_5 &= &0  \\
          &   &2c_2  &+  &2c_3  &-  &c_4 &+  &3c_5 &= &0  \\
          &   &      &   &     &   &c_4  &   &     &= &0  
\end{linsys}
\end{equation*}
其解集具有这一参数化。
\begin{equation*}
  \set{\colvec{c_1 \\ c_2 \\ c_3 \\ c_4 \\ c_5}=
     c_3\colvec[r]{-1 \\ -1 \\ 1 \\ 0 \\ 0}
     +c_5\colvec[r]{-3 \\ -3/2 \\ 0 \\ 0 \\ 1}
     \suchthat c_3,c_5\in\Re }
\end{equation*}
取 $c_5=1$ 和~$c_3=0$，
得到 ($*$) 的一个实例。
\begin{equation*}
  -3\cdot\colvec[r]{1 \\ 0 \\ 0}
  -\frac{3}{2}\cdot\colvec[r]{0 \\ 2 \\ 0}
  +0\cdot\colvec[r]{1 \\ 2 \\ 0}
  +0\cdot\colvec[r]{0 \\ -1 \\ 1}
  +1\cdot\colvec[r]{3 \\ 3 \\ 0}
  =\colvec[r]{0 \\ 0 \\ 0}
\end{equation*}
这表明 $S$ 中我们与~$c_5$
相关联的
那个向量，位于由 $c_1$ 的向量与 $c_2$ 的向量组成的集合的张成之中。
我们可以舍弃 $S$ 的第五个向量而不使张成缩小。

类似地，取 $c_3=1$，$c_5=0$，
得到~($*$) 的一个实例，
它表明我们可以舍弃 $S$ 的第三个向量而不使张成缩小。
 于是这个集合与 $S$ 有相同的张成。
\begin{equation*}
  % \hat{S}=
     \set{\colvec[r]{1 \\ 0 \\ 0},
     \colvec[r]{0 \\ 2 \\ 0},
     \colvec[r]{0 \\ -1 \\ 1}  }
\end{equation*}
验证它是线性无关的是常规的。
\end{example}

\begin{corollary} \label{cor:LDMeansLC}
%<*co:LDMeansLC>
向量空间的子集 \( S=\set{\vec{s}_1,\dots,\vec{s}_n} \)
是线性相关的，当且仅当某个 \( \vec{s_i} \)
是排在其之前的那些向量
\( \vec{s}_1 \), \ldots, \( \vec{s}_{i-1} \)
的线性组合。
%</co:LDMeansLC>
\end{corollary}

\begin{proof}
%<*pf:LDMeansLC>
考虑 \( S_0=\set{} \)，\( S_1=\set{\vec{s_1}} \)，
\( S_2=\set{\vec{s}_1,\vec{s}_2 } \)，等等。
某个下标 \( i\geq 1 \) 是第一个使
\( S_{i-1}\union\set{\vec{s}_i } \)
线性相关的，此时 \( \vec{s}_i\in\spanof{ S_{i-1} } \)。
%</pf:LDMeansLC>
\end{proof}

\nearbycorollary{th:AlwaysAnLDSubset} 的证明描述了
通过缩小、通过取子集来产生一个线性
无关集。
而 \nearbycorollary{cor:LDMeansLC} 的证明描述了通过取扩集
来找出一个线性相关集。
我们以如下问题结束本小节：
线性无关与线性相关如何相互作用
与集合之间的子集关系。

\begin{lemma}  \label{le:SubsetPreserveLI}
%<*lm:SubsetPreserveLI>
线性无关集的任何子集也是线性无关的。
线性相关集的任何扩集也是线性相关的。
%</lm:SubsetPreserveLI>
\end{lemma}

\begin{proof}
%<*pf:SubsetPreserveLI>
两者都是显然的。
%</pf:SubsetPreserveLI>
\end{proof}

%<*SubsetPreserveLI>
重述一下：子集保持无关性，
扩集保持相关性。
%</SubsetPreserveLI>

这是四种可能情形中的两种。
第三种情形，子集是否保持线性相关，
由 \nearbyexample{ex:ShrinkSetSameSpan} 覆盖，它给出了
一个线性相关集 $S$，它的一个子集
是线性相关的，而另一个
是无关的。
第四种情形，扩集是否保持线性无关，
由 \nearbyexample{ex:LinindSetsAndSuper} 覆盖，它给出了一些例子：
一个线性无关集既有无关的扩集，
又有相关的扩集。
下表对此进行了总结。
\smallskip
%<*IndependenceAndSubsetTable>
\begin{center} \renewcommand{\arraystretch}{1.1}
  \begin{tabular}[b]{r|cc} 
                        \multicolumn{1}{c}{}
                        &\multicolumn{1}{c}{\( \hat{S}\subset S \)}
                        &\multicolumn{1}{c}{\( \hat{S}\supset S \)}      \\
     \cline{2-3}\rule{0em}{12pt} % make box a bit taller
          \textit{$S$ 无关} 
              &\( \hat{S} \) 必为无关   &\( \hat{S} \) 两者皆可\\
     % \cline{2-3}\noalign{\smallskip}
          \textit{$S$ 相关} 
              &\( \hat{S} \) 两者皆可 &\( \hat{S} \) 必为相关    \\
     % \cline{2-3}
   \end{tabular}
\end{center}
%</IndependenceAndSubsetTable>
\smallskip

关于线性无关与扩集之间的相互作用，
\nearbyexample{ex:LinindSetsAndSuper} 还说明了些什么。
它指出了这样一个
线性无关集，它之所以是极大的，
是因为它没有线性无关的扩集。
由 \nearbylemma{lm:AddVecLiSetIsLiIffVecNotInSpan}，
一个线性无关集是极大的，当且仅当它
张成
整个空间，因为只有这时，
空间中的所有向量才都已经在这个张成之中。
这很好地
补充了 \nearbylemma{cor:LiSetMinSpanSet}，即
张成集是最小的当且仅当它是线性无关的。

% In summary,
% we have introduced the definition of linear independence to 
% formalize the idea of the minimality of a spanning set.
% We have developed some properties of this idea.
% The most important is \nearbylemma{lm:AddVecLiSetIsLiIffVecNotInSpan}, which 
% tells us that a linearly independent set is maximal when it spans the space.

